\section{Un programma di verifica}\label{sec:future-work}
\subsection{Dal modello al luogo reale}
Il modello Blender va trattato come rappresentazione progettuale
versionata. Finché mancano quote affidabili, le proporzioni permettono
confronti tra ipotesi, non attestazioni dimensionali. Evitiamo quindi
di chiamarlo \emph{digital twin} senza una definizione e una calibrazione
adeguate. Un render non misura riverbero, temperatura o accessibilità.

Il prossimo passaggio è associare a ogni geometria la sua provenienza:
rilievo, fotografia, stima o elemento provvisorio. Le alternative di
collocazione potranno poi essere confrontate sulla stessa base, conservando
la storia delle correzioni. Le verifiche acustiche e impiantistiche
richiederanno dati e competenze specifiche.

\subsection{Un primo prototipo circoscritto}
Una prova temporanea e presidiata è una possibilità da valutare prima
di un'installazione stabile. Permetterebbe di osservare difficoltà
d'ingresso, gestione dell'audio, uso a due e necessità di assistenza.
Non viene qui presentata come attività già autorizzata o programmata.

La selezione dei partecipanti dovrebbe includere diversi livelli di
familiarità musicale e digitale. Le esigenze di accesso vanno considerate
nella progettazione delle prove, evitando di assumere che un'interfaccia
semplice per il progettista lo sia anche per chi la incontra per la prima volta.
Prima della prova andranno definite le modalità di partecipazione e
quali dati sia davvero necessario raccogliere.

\subsection{Che cosa osservare e che cosa non confondere}
Per l'usabilità proponiamo di registrare, con un protocollo proporzionato,
l'avvio della prima attività, le richieste d'aiuto e gli abbandoni.
Per la collaborazione interessa osservare scambio dei ruoli e aiuto
reciproco. Per l'apprendimento occorre invece scegliere prima un obiettivo
musicale circoscritto e verificare se la persona riesce a svolgerlo anche
senza la guida della postazione. Un'eventuale verifica successiva
permetterebbe di distinguere l'esecuzione immediata da ciò che rimane.

Una prima prova potrà documentare fattibilità e problemi d'uso.
Per sostenere un effetto specifico del Portale serviranno invece un
disegno comparativo adeguato e un controllo delle spiegazioni alternative:
assistenza ricevuta, esperienza precedente, novità della situazione.
Numerosità, misure e criteri decisionali dovranno essere stabiliti
prima della raccolta, in funzione della domanda scelta.

La sostenibilità richiede un'altra serie di dati: tempo di presidio,
guasti, aggiornamenti dei materiali, costi e compatibilità con gli altri
usi del parco. Un prototipo piacevole ma impossibile da mantenere
non costituisce ancora un servizio civico.

\subsection{Passaggi necessari}
Il programma propone tre passaggi, ciascuno con un risultato controllabile:
\begin{enumerate}
\item Consolidare le fonti: identificare documenti e pagine per cronologia,
proprietà e tutela; aggiornare il registro dei claim.
\item Verificare il luogo: raccogliere quote e condizioni d'uso, confrontare
le collocazioni e motivare l'eventuale esclusione del Padiglione.
\item Definire una prova: delimitare attività, responsabilità, materiali,
costi e protocollo di osservazione prima dell'installazione.
\end{enumerate}

\paragraph{Stato editoriale e contributo degli strumenti.}
Questa è la prima bozza italiana. La scrittura e il controllo dei riferimenti
sono assistiti da ChatGPT; Pi interviene nella revisione e nella verifica
tecnica. Questo confronto non equivale a una peer review scientifica
indipendente. Le responsabilità delle affermazioni rimangono dell'autore.

Il lavoro prosegue con il controllo documentale delle lacune storiche
e con la raccolta dei dati necessari a discutere una prova sul posto.
La scelta del Padiglione potrà essere confermata, modificata o abbandonata
alla luce di quei dati.
